% ============================================================================
% Fichier  : partie3/P03_C07_identifier_preserver.tex : FORN (LPI342)
% Sources : ISO/IEC 27037 ; RFC 3227 ; ACPO. Opération TSANAGA.
% ============================================================================

\chapter{Identifier et Préserver les Sources}
\label{chap:preserver}
\chaptermark{Identifier et préserver}

\epigraph{%
    « La preuve la plus solide est celle que l'on n'a pas abîmée\\
    en voulant la regarder. »%
}{Principe du cours}

\niveauChapitre{Fondamental}{Partie~II.}

\begin{boxobjectifs}
\begin{itemize}
    \item Identifier les sources potentielles de preuve.
    \item Préserver une scène numérique sans l'altérer.
    \item Décider d'éteindre, de débrancher ou de capturer à chaud.
    \item Protéger contre l'écriture (bloqueur).
    \item Documenter l'état initial.
\end{itemize}
\end{boxobjectifs}

\bigskip

% ============================================================================
\section*{Question directrice : devant l'ordinateur de M. NDONGO, faut-il l'éteindre, le débrancher, ou le laisser allumé ?}
% ============================================================================

Chaque geste sur un système vivant a un coût. Ce chapitre apprend à
identifier les sources, à préserver la scène et à décider de la conduite à
tenir selon l'état du support.

\medskip
\textbf{Plan du chapitre.}
\begin{enumerate}
    \item Les sources de preuve (section~\ref{sec:c07-sources}).
    \item Préserver la scène (section~\ref{sec:c07-scene}).
    \item Éteindre, débrancher ou capturer (section~\ref{sec:c07-decision}).
    \item Protéger contre l'écriture (section~\ref{sec:c07-bloqueur}).
\end{enumerate}

\separateur

% ============================================================================
\section{Les Sources de Preuve}
\label{sec:c07-sources}
% ============================================================================

\begin{tableaucours}{L{3.6cm}L{8.6cm}}{%
    \tableauHeader{Source} &
    \tableauHeader{Ce qu'elle peut contenir}
}
Disques et clés USB & Fichiers, fragments, artefacts, fichiers supprimés \\
Mémoire vive & Processus, secrets, données non écrites \\
Téléphones & Messages, applications, géolocalisation \\
Équipements réseau & Journaux, flux, configurations \\
Services en ligne & Comptes, journaux, sauvegardes \\
\end{tableaucours}
\vspace{-0.8em}
\begin{center}{\small\itshape Tableau~7.1. Sources potentielles de preuve}\end{center}

% ============================================================================
\section{Préserver la Scène}
\label{sec:c07-scene}
% ============================================================================

\begin{boxprocedure}
\textbf{Préserver une scène numérique}
\begin{enumerate}
    \item Photographier l'état initial (écran, branchements, voyants).
    \item Ne toucher à rien avant d'avoir décidé et documenté.
    \item Isoler des réseaux (éviter effacement ou verrouillage à distance).
    \item Noter l'heure du système et l'heure réelle (écart).
    \item Consigner chaque geste.
\end{enumerate}
\end{boxprocedure}

% ============================================================================
\section{Éteindre, Débrancher ou Capturer}
\label{sec:c07-decision}
% ============================================================================

\begin{tableaucours}{L{3.2cm}L{4.6cm}L{4.4cm}}{%
    \tableauHeader{État} &
    \tableauHeader{Conduite} &
    \tableauHeader{Raison}
}
Éteint & Ne pas allumer, acquérir le disque & L'allumage modifie le disque \\
Allumé, données volatiles utiles & Capturer la mémoire d'abord & La mémoire meurt à l'extinction \\
Allumé, risque de chiffrement & Éviter l'extinction brutale & Peut déclencher un verrouillage \\
\end{tableaucours}
\vspace{-0.8em}
\begin{center}{\small\itshape Tableau~7.2. Décider selon l'état du support}\end{center}

\begin{boxalert}
\textbf{Violation typique.} Allumer l'ordinateur éteint « juste pour voir ce
qu'il contient ».

\smallskip
\textbf{Règle~:} allumer modifie des dizaines d'artefacts (dates d'accès,
démarrage). On n'allume pas un support éteint~: on acquiert son disque. On ne
touche au vivant que pour une raison documentée.
\end{boxalert}

% ============================================================================
\section{Protéger contre l'Écriture}
\label{sec:c07-bloqueur}
% ============================================================================

\begin{boxscenario}{Opération TSANAGA. Le bloqueur en écriture}
\textbf{Constatation.} Pour lire le disque et la clé saisis, Mme~DJUIDJE passe
par un \termecle{bloqueur en écriture} (matériel ou logiciel), qui autorise la
lecture mais interdit toute modification du support.

\medskip
\textbf{Analyse} (confiance élevée). Sans bloqueur, le simple branchement
aurait pu écrire sur le support (horodatage, index). Le bloqueur garantit
l'intégrité exigée pour la recevabilité (Chapitre~\ref{chap:recevabilite}).
\end{boxscenario}

\begin{boxmetier}{Analyste forensique}
\textbf{Lire sans écrire.} L'analyste branche toujours un support de preuve
derrière un bloqueur en écriture, et vérifie par condensat qu'il n'a rien
modifié. C'est un réflexe, pas une option.
\end{boxmetier}

\separateur

\begin{boxresume}
\textbf{Ce qu'il faut retenir de ce chapitre~:}
\begin{itemize}
    \item \textbf{Les sources}~: disques, mémoire, téléphones, réseau,
          services.
    \item \textbf{Préserver la scène}~: photographier, isoler, ne rien
          toucher sans raison.
    \item \textbf{Éteint}~: ne pas allumer. \textbf{Allumé}~: capturer la
          mémoire d'abord.
    \item \textbf{Un bloqueur en écriture} protège l'intégrité.
    \item \textbf{On documente} l'état initial et chaque geste.
\end{itemize}
\end{boxresume}

\begin{boxinfo}
\textbf{Points de contrôle de l'élève.} Je sais~: citer les sources~;
préserver une scène~; décider selon l'état~; utiliser un bloqueur~;
documenter l'état initial.
\end{boxinfo}

\bigskip

\begin{boxexo}[title={\ding{45}\quad Exercice~7.1. Sources \niveaubase}]
Citez quatre sources potentielles de preuve.
\end{boxexo}

\begin{boxexo}[title={\ding{45}\quad Exercice~7.2. Éteint \niveaubase}]
Que faire d'un ordinateur éteint~? Pourquoi~?
\end{boxexo}

\begin{boxexo}[title={\ding{45}\quad Exercice~7.3. Bloqueur \niveaubase}]
À quoi sert un bloqueur en écriture~?
\end{boxexo}

\begin{boxexo}[title={\ding{45}\quad Exercice~7.4. Allumé \niveaumoyen}]
Un ordinateur est allumé avec des données volatiles utiles. Que faire~?
\end{boxexo}

\begin{boxexo}[title={\ding{45}\quad Exercice~7.5. Allumer \niveaumoyen}]
Pourquoi n'allume-t-on pas un support éteint pour « voir »~?
\end{boxexo}

\begin{boxexo}[title={\ding{45}\quad Exercice~7.6. Isoler \niveaumoyen}]
Pourquoi isoler un support des réseaux dès la saisie~?
\end{boxexo}

\begin{boxexo}[title={\ding{45}\quad Exercice~7.7. Chiffrement \niveaumoyen}]
Pourquoi éviter l'extinction brutale d'un système peut-être chiffré~?
\end{boxexo}

\begin{boxexo}[title={\ding{45}\quad Exercice~7.8. État initial \niveaumoyen}]
Que consigne-t-on de l'état initial d'une scène~?
\end{boxexo}

\begin{boxexo}[title={\ding{45}\quad Exercice~7.9. Procédure de saisie \niveauexpert}]
Rédigez la procédure de saisie et de préservation pour TSANAGA.
\end{boxexo}

\begin{boxexo}[title={\ding{45}\quad Exercice~7.10. Fiche de décision \niveauexpert}]
Concevez une fiche d'aide à la décision « éteint / allumé / chiffré ».
\end{boxexo}

\bigskip

\begin{boxinfo}
\textbf{Pour aller plus loin.}
\begin{itemize}
    \item ISO/IEC 27037~\cite{iso27037}~; RFC~3227~\cite{rfc3227}~;
          ACPO~\cite{acpo2012}.
\end{itemize}
\end{boxinfo}

\begin{boxinfo}
\textbf{Transition.} Une fois le support préservé, on le duplique~: le
Chapitre~\ref{chap:duplication} traite de la duplication et du hachage.
\end{boxinfo}
